\caption{\textbf{A local inhibitory-context gate supplies selective credit in a conductance tree.}
\textbf{A}, Four terminal compartments feed two proximal compartments and a soma. The hard rule sends the somatic error only to terminals below the uninhibited parent. Both proximal compartments retain unit credit multipliers, so all 24 conductances remain plastic. The illustrated context switches between the two branches.
\textbf{B}, Nominal teacher tuning. Terminal E/I conductances produce aligned or opposed feature responses; both subtrees receive identical positive input features.
\textbf{C,D}, Mean held-out NMSE during Adam training at rate 0.03 in twenty new paired seed blocks; shading gives descriptive 95\% seed-bootstrap intervals. The dotted line marks 4,096 updates; every condition continues to 16,384. Two weighted leaf profiles use oracle coefficients projected from the exact field and unit proximal credit. The local gate evaluates no exact paths during learning.
\textbf{E}, Paired task-by-rule interaction: broadcast-minus-local-gate NMSE on opposed minus aligned targets. Checkpoints minimize validation error within each budget. Dots are twenty paired seeds; diamonds and whiskers show means and descriptive 95\% seed-bootstrap intervals.
\textbf{F}, Opposed-task endpoints within 4,096 updates. A swapped selector credits the inhibited subtree. Applying the gate to proximal credit also suppresses the two inhibitory-gain gradients by construction. All outcomes are retained. Markers and intervals follow \textbf{E}. The experiment tests a supplied local context signal, not endogenous discovery of that signal. Nine-rule, rate and longer-budget outcomes are supplied in Supplementary Fig.~\ref{fig:supp_local_gate_rates}.}
